% =====================================================================
\section{How Our ROS~2 Node Actually Controls the Drone}
\label{sec:realctrl}
% =====================================================================

\Cref{sec:model1} explained each equation. This section follows our real program, \code{drones\_controller},
from start-up through one complete control step, line by line, with numbers. At the end you should be able to point
at any line of \code{controller\_node.py} or \code{so3\_controller.py} and say which equation it is and why it is there.

\subsection{Start-up: reading the parameters and building the controller}

When the node starts, it declares every parameter with a default value. ROS~2 then replaces the defaults with the
values in \code{controller.yaml}. For example:
\pyfile{94}{122}{\node}{controller\_node.py --- some of the parameter declarations}
Then it creates the two objects that do the work: the \emph{trajectory generator} (where should I be?) and the
\emph{SO(3) controller} (how do I get there?). The gains are turned into diagonal matrices $K_p,K_v,K_R,K_\omega$ and the
inertia into $J$:
\pyfile{242}{255}{\node}{controller\_node.py --- the two working objects}
\pyfile{19}{30}{\so}{so3\_controller.py --- \_\_init\_\_(): gains become diagonal matrices}
Next it subscribes to the estimator's output and starts a timer that calls \code{control\_loop()} at
\code{control\_rate} = 25 times per second (one call every 40\,ms):
\pyfile{261}{273}{\node}{controller\_node.py --- subscribing to the state estimate}
\pyfile{338}{343}{\node}{controller\_node.py --- the control timer}

\begin{warn}[Small inconsistency to know about]
The default \code{control\_rate} in the code is 50\,Hz, but \code{controller.yaml} sets 25\,Hz. When the node is launched
with the YAML file the rate is 25\,Hz. If it is started without the YAML, it runs at 50\,Hz. The YAML value is the one
we use everywhere in this document.
\end{warn}

\subsection{Between control steps: receiving the state (block 4's output)}

Whenever ArduPilot publishes a new estimate, a \emph{callback} stores it. \code{pose\_callback} stores $\vp$ and turns the
quaternion into $R$ (shown in \cref{sec:m1b1}). \code{twist\_callback} stores $\vv$ and $\vw$:
\pyfile{419}{433}{\node}{controller\_node.py --- twist\_callback()}
Both remember \emph{when} the data arrived, so that the control step can refuse to use old data.

\subsection{One control step, line by line}

Every 40\,ms the timer calls \code{control\_loop()}:
\pyfile{526}{543}{\node}{controller\_node.py --- control\_loop(), first part}
\begin{enumerate}
  \item \code{state\_is\_valid()}: is every part of the state present, younger than 0.5\,s and a finite number? If not,
        \emph{do nothing} this step (safety).
  \item \code{get\_desired\_state()}: the mission. Today it is a fixed waypoint, so $\vp_d=[0,0,2]$, $\vv_d=\va_d=\bm0$,
        $\psi_d=0$, $\vw_d=\dot\vw_d=\bm0$:
        \pyfile{11}{29}{\traj}{trajectory\_generator.py --- get\_desired\_state()}
  \item \code{controller.compute()}: blocks 5, 6 and 7, in that order:
        \pyfile{221}{258}{\so}{so3\_controller.py --- compute(): blocks 5 $\rightarrow$ 6 $\rightarrow$ 7}
  \item Publish every intermediate result on \code{/drones\_controller/...}, so we can plot them while flying.
  \item If \code{enable\_command} is true, send a velocity command (next subsection).
\end{enumerate}

\begin{toy}[one complete control step of our node, with real numbers]
The estimate says the drone is at $\vp=[0.1,\,-0.2,\,1.5]$\,m, climbing at $\vv=[0,0,0.3]$\,m/s, level ($R=I$) and not
turning ($\vw=\bm0$). The waypoint is $[0,0,2]$. Running the real \code{SO3Controller.compute()} on these numbers gives:
\begin{center}\small
\begin{tabular}{@{}l l l l@{}}
\toprule
\textbf{Output} & \textbf{Equation} & \textbf{Value} & \textbf{Meaning}\\
\midrule
\code{e\_p} & (15) & $[-0.1,\ 0.2,\ 0.5]$\,m & target is 10\,cm west, 20\,cm north, 50\,cm up\\
\code{e\_v} & (15) & $[0,\ 0,\ -0.3]$\,m/s & already climbing a bit\\
\code{F\_d} & (16)--(17) & $[-0.225,\ 0.450,\ 15.540]$\,N & push mostly up, slightly towards the target\\
\code{R\_d} & (18)--(19) & tilt of 1.85$^\circ$ & lean slightly towards $-x$, $+y$\\
\code{e\_R} & (20) & $[0.0289,\ 0.0145,\ -0.0002]$ & the drone is level but should lean\\
\code{e\_omega} & (21) & $[0,\ 0,\ 0]$ & no spin, none wanted\\
\code{M\_d} & (22) & $[-0.1158,\ -0.0579,\ 0.0004]$\,N\,m & twist towards the lean\\
\bottomrule
\end{tabular}
\end{center}
Check the torque: $-K_R\ve_R=-[4(0.0289),\ 4(0.0145),\ 2(-0.0002)]=[-0.1158,\ -0.0579,\ 0.0004]$, and $\ve_\omega=0$ adds
nothing. A torque about $-x$ tips the top of the drone towards $+y$, and a torque about $-y$ tips it towards $-x$: exactly the
direction of $\vF_d$. Every number agrees with the toy model (\cref{sec:m1b5,sec:m1b6,sec:m1b7}).
\end{toy}

\subsection{What is actually sent to the drone: the velocity command}

ArduPilot's ROS~2 interface \code{/ap/v1/cmd\_vel} accepts a \emph{velocity}, not a force and torque. So our node does
not send $\vF_d$ or $\vtau$ (they are published only for monitoring). Instead \code{publish\_velocity\_command()} turns
the position and velocity errors into a velocity request:
\pyfile{670}{696}{\node}{controller\_node.py --- publish\_velocity\_command(): the velocity law}
\[
\vv_{\text{cmd}}=0.8\,\ve_p+0.2\,\ve_v,\qquad\text{and if }\|\vv_{\text{cmd}}\|>1\text{ m/s, scale it down to length 1 m/s}.
\]
This is again a spring (0.8: ``go towards the target, faster when further away'') and a damper (0.2: ``slow down if you are
already moving''). The scaling keeps the \emph{direction} and enforces \code{max\_velocity}.

\begin{toy}[the velocity command]
\textbf{Close to the target} (same state as above): $\vv_{\text{cmd}}=0.8[-0.1,0.2,0.5]+0.2[0,0,-0.3]=[-0.08,\ 0.16,\ 0.34]$\,m/s,
length 0.38\,m/s, below the limit, sent as it is.

\textbf{Far from the target}: drone at $[3,\,-2,\,0.5]$, at rest. $\ve_p=[-3,\,2,\,1.5]$, so
$\vv_{\text{cmd}}=[-2.4,\ 1.6,\ 1.2]$, length 3.12\,m/s, too fast. Scaled by $1/3.12$: $[-0.768,\ 0.512,\ 0.384]$\,m/s,
length exactly 1\,m/s, still pointing straight at the target.
\end{toy}

ArduPilot then does the rest. Inside the flight controller it runs its own version of blocks 5--8: a velocity controller
that asks for an acceleration, which becomes a lean angle, which an attitude controller turns into torques, which the
mixer turns into motor speeds. So the whole chain of \cref{sec:model1} still happens; our node does the outer
``where should I go'' part and ArduPilot the inner, fast part.

\begin{warn}[Two things to check in the simulator]
\begin{itemize}
  \item \textbf{The sign of $z$.} The last lines send \code{linear.z = -velocity\_command[2]}:
        \pyfile{709}{719}{\node}{controller\_node.py --- filling the message}
        The pose and twist that we read are in ENU ($z$ up), and the header says \code{"map"}. If ArduPilot reads this
        command in ENU too, the minus sign makes the drone go \emph{down} when it should go up. It is only right if the
        interface expects $z$ pointing down (NED). This must be checked once in SITL: command a small climb and watch
        the altitude.
  \item \textbf{Safety switch.} \code{enable\_command} is \code{false} in the YAML, so by default the node only computes and
        publishes; nothing moves until we deliberately switch it on.
\end{itemize}
\end{warn}

\subsection{Summary: from equation to line of code}

\begin{center}\small
\begin{tabularx}{\linewidth}{@{}l l X@{}}
\toprule
\textbf{Equation} & \textbf{Block} & \textbf{Where in our code}\\
\midrule
(1)--(4) & dynamics & Gazebo (not in our code); result read in \code{pose\_callback}, \code{twist\_callback}\\
$R$ from quaternion & -- & \code{state\_adapter.py}: \code{quaternion\_to\_rotation\_matrix()}, lines 4--42\\
(5)--(14) & EKF (VIO) & ArduPilot EKF3; freshness check in \code{state\_is\_valid()}, lines 447--496\\
(15) & position errors & \code{so3\_controller.py}: \code{compute\_position\_errors()}, lines 32--40\\
(16)--(17) & desired force & \code{so3\_controller.py}: \code{compute\_desired\_force()}, lines 42--74\\
(18)--(19) & desired attitude & \code{so3\_controller.py}: \code{compute\_desired\_attitude()}, lines 76--126\\
(20) & attitude error & \code{so3\_controller.py}: \code{compute\_attitude\_error()}, lines 128--142\\
(21) & rate error & \code{so3\_controller.py}: \code{compute\_angular\_velocity\_error()}, lines 144--158\\
(22) & torque & \code{so3\_controller.py}: \code{compute\_moment()}, lines 160--210\\
(23)--(25) & thrust, mixer & ArduPilot (our node sends a velocity instead)\\
$\vv_{\text{cmd}}$ & command & \code{controller\_node.py}: \code{publish\_velocity\_command()}, lines 644--734\\
\bottomrule
\end{tabularx}
\end{center}
